% Gerado por ferramentas/importar_relatorios.py a partir de MLP/artigo/MLP_EA_bomba.tex.
% Revisado à mão na dissertação (versão enxuta): edite este arquivo, não a fonte.
\subsection{Dados e delineamento}
\label{mlp:sec:dados}

A bancada \completar{descrever: bomba, motor e inversor, tubulação, sensores de EA, posição dos canais A e B, pré-amplificadores} foi operada em delineamento cruzado: oito frequências do motor (20 a 90\,Hz, em passos de 10\,Hz) combinadas com sete configurações de abertura das válvulas de recalque e sucção. Em cada configuração, uma válvula ficava totalmente aberta e a outra em 75, 50 ou 25\,\%; no formato recalque-sucção, as configurações são 100-100, 100-75, 100-50, 100-25, 75-100, 50-100 e 25-100. Cada combinação foi registrada uma única vez, por cerca de 300\,s, em 56 corridas. O MOMEA registrou dois canais (A e B), com ganho de 34\,dB, HDT de 250\,\(\mu\)s, HLT re-armável de 2\,ms e MAXDUR de 100\,ms; os rótulos foram lidos da estrutura de pastas.

Cada série de válvulas foi ensaiada em bloco contínuo, em dois dias (Tabela~\ref{mlp:tab:ordem}), com as frequências em ordem crescente ou decrescente. Assim, as válvulas estão confundidas com o dia, o horário e o limiar de detecção, e a frequência está confundida com o tempo dentro de cada série, mas não entre séries, porque o sentido da varredura muda.

\begin{table}[htbp]
	\centering
	\caption{Séries de válvulas, horários, sentido da varredura de frequência e limiares.}
	\label{mlp:tab:ordem}
	\small
	\begin{tabular}{llll}
		\toprule
		\textbf{Série} & \textbf{Dia e horário} & \textbf{Ordem das freq.} & \textbf{Limiar A/B (dB\(_{\mathrm{AE}}\))} \\
		\midrule
		100-100 & 05/11, 10h42--12h09 & 30\(\rightarrow\)90, depois 20 & 63/50 (20\,Hz: 65/55) \\
		25-100 & 05/11, 15h54--16h41 & 90\(\rightarrow\)20 & 65/55 \\
		100-50 & 05/11, 16h52--17h45 & 90\(\rightarrow\)20 & 65/55 \\
		100-25 & 06/11, 09h49--10h32 & 20\(\rightarrow\)90 & 50/50 \\
		100-75 & 06/11, 11h17--12h01 & 20\(\rightarrow\)90 & 50/50 \\
		75-100 & 06/11, 14h49--15h32 & 20\(\rightarrow\)90 & 50/50 \\
		50-100 & 06/11, 15h52--16h41 & 20\(\rightarrow\)90 & 50/50 \\
		\bottomrule
	\end{tabular}
\end{table}

\subsection{Qualidade dos dados}
\label{mlp:sec:qualidade}

Cinco características dos registros condicionam a análise:
\begin{itemize}
	\item \textbf{saturação}: 69,8\,\% dos \textit{hits} têm o indicador SAT; a partir de 40 a 50\,Hz, praticamente todos saturam e a amplitude acumula-se no fundo de escala, 102,8\,dB\(_{\mathrm{AE}}\) (Figuras~\ref{mlp:fig:exemplo} e~\ref{mlp:fig:amplitude}), deixando de discriminar as condições;
	\item \textbf{duração máxima}: 17,8\,\% dos \textit{hits}, quase todos no canal A, foram encerrados pelo MAXDUR, o que indica emissão contínua, coerente com ruído de escoamento ou cavitação;
	\item \textbf{rearme}: o ROVL está ativo em todos os \textit{hits}, e os \textit{hits} chegam espaçados em cerca de 0,1\,s, de modo que o detector passa a amostrar um sinal quase contínuo a cerca de 10\,Hz e a taxa de \textit{hits} deixa de medir a atividade acústica;
	\item \textbf{limiar variável}: o limiar mudou entre sessões (Tabela~\ref{mlp:tab:ordem}); como a duração, as contagens e a energia são medidas acima do limiar da sessão, esse efeito não pode ser corrigido depois;
	\item \textbf{tempo}: o carimbo de tempo está zerado, e usou-se a hora de chegada no computador, suficiente para janelas de segundos.
\end{itemize}
A Figura~\ref{mlp:fig:mapas} resume esses indicadores por condição.

\begin{figure}[htbp]
	\centering
	\includegraphics[width=\textwidth]{figs/mlp/fig_exemplo_hits.png}
	\caption{Trechos de 10\,s de \textit{hits} em duas corridas. A 20\,Hz (esquerda), as amplitudes variam e não há saturação. A 80\,Hz (direita), todos os \textit{hits} estão no fundo de escala e quase todos os do canal A atingem a duração máxima (pontos pretos). As linhas pontilhadas marcam as janelas de 5\,s.}
	\label{mlp:fig:exemplo}
\end{figure}

\begin{figure}[htbp]
	\centering
	\includegraphics[width=0.9\textwidth]{figs/mlp/fig_amplitude_saturacao.png}
	\caption{Distribuição da amplitude de pico por faixa de frequência do motor. A partir de 40\,Hz, a maior parte dos \textit{hits} acumula-se em 102,8\,dB\(_{\mathrm{AE}}\) (saturação do ADC).}
	\label{mlp:fig:amplitude}
\end{figure}

\begin{figure}[htbp]
	\centering
	\includegraphics[width=\textwidth]{figs/mlp/fig_mapas_condicao.png}
	\caption{Indicadores por condição, um valor por corrida: taxa de \textit{hits}, fração de \textit{hits} saturados (SAT), fração com duração máxima (MAXDUR) e mediana de \(\log_{10}\) da energia por segundo no canal A.}
	\label{mlp:fig:mapas}
\end{figure}

\subsection{Descritores por janela}
\label{mlp:sec:descritores}

Como a cavitação e o ruído de escoamento são quase contínuos, um \textit{hit} isolado não carrega a condição. Os \textit{hits} foram agregados em janelas consecutivas de \(T_w = 5\)\,s, sem sobreposição, o que resulta em 3\,540 janelas (cerca de 63 por corrida, com 25 a 50 \textit{hits} por canal). Para cada janela e canal calcularam-se 18 descritores, 36 no total:
\begin{itemize}
	\item \textbf{com todos os \textit{hits}}: taxa de \textit{hits}, frações SAT e MAXDUR, mediana da frequência média \(\mathrm{FM} = N_{\mathrm{cont}}/t_{\mathrm{dur}}\) e o logaritmo da energia por segundo,
	\begin{equation}
		E_s = \log_{10}\!\Big(\frac{1}{T_w}\sum_{i\in w} E_i + 1\Big),
		\label{mlp:eq:energia_s}
	\end{equation}
	que funciona como um indicador de nível RMS;
	\item \textbf{só com os \textit{hits} sem MAXDUR}, cuja duração não foi truncada: média e percentis 10, 50 e 90 da amplitude; percentis 50 e 90 de \(\log_{10}\) da energia, da duração e das contagens; medianas de \(\log_{10}\) do tempo de subida, da frequência média e de \(\log_{10}\) de \(\mathrm{RA} = t_{\mathrm{sub}}/A_{\mathrm{pico}}\) \cite{ohno2010}.
\end{itemize}
Janelas sem \textit{hits} normais (6,7\,\% dos valores) têm descritores imputados pela mediana; imputação e padronização são ajustadas apenas no treino de cada partição. A energia por segundo cresce de forma monotônica com a frequência e estabiliza a partir de cerca de 60\,Hz; por válvula, as diferenças são menores que a dispersão entre frequências.

\subsection{Alvos, modelos e validação}
\label{mlp:sec:metodo}

A \textbf{frequência} do motor, grandeza ordenada, foi tratada como regressão, com leitura ordinal pelo arredondamento ao nível mais próximo, de modo que errar 40 por 50\,Hz custe menos que errar por 90\,Hz. As \textbf{válvulas} foram modeladas como uma classe de 7 níveis, equilibrada com 8 corridas por classe, porque o delineamento cruzado só contém 7 das 16 combinações possíveis; os acertos em \% recalque e \% sucção foram derivados da classe prevista.

Foram comparados três modelos com hiperparâmetros fixos, precedidos de imputação e padronização no mesmo \textit{pipeline} \cite{pedregosa2011}: uma \textbf{MLP} com duas camadas ocultas de 64 e 32 neurônios, ReLU, Adam, penalidade L2 \(\alpha = 10^{-3}\) e parada antecipada com 15\,\% do treino e paciência de 30 épocas \cite{prechelt1998}; uma \textit{Random Forest} (RF) de 300 árvores \cite{breiman2001}; e um \textit{gradient boosting} por histogramas (GB) \cite{friedman2001}. Não houve busca de hiperparâmetros: com 56 corridas, uma busca aninhada teria pouca resolução \cite{varma2006}, e os números são linhas de base.

Nenhuma validação sorteia \textit{hits} ou janelas entre treino e teste. Usaram-se três esquemas, em ordem crescente de independência \cite{roberts2017}:
\begin{description}
	\item[BLOCOS (5 partições).] Cada corrida é dividida em 5 trechos consecutivos de cerca de 60\,s; na partição \(k\), o \(k\)-ésimo trecho de todas as corridas é o teste, com uma janela de exclusão de cada lado. A mesma corrida aparece no treino e no teste.
	\item[LORO (\textit{leave-one-run-out}, 56 partições).] Uma corrida inteira é o teste; sua frequência aparece no treino com outras válvulas, e sua configuração, com outras frequências. É o teste de generalização para uma condição não vista.
	\item[LOSO (\textit{leave-one-series-out}, 7 partições, só frequência).] Uma série de válvulas inteira, isto é, uma sessão contínua, é o teste.
\end{description}
Para a frequência, as métricas são MAE, RMSE, \(R^2\), acerto do nível exato e acerto em \(\pm 10\)\,Hz; para as válvulas, acurácia, acurácia balanceada, F1 macro e acertos de \% recalque e \% sucção. Todas foram calculadas por janela e por corrida (mediana das previsões ou voto majoritário).
